\centerline{\bf \Huge Степень точки относительно...точки?}
\centerline{\bf \Huge }

\begin{flushright}
        \scriptsize{}
\end{flushright}


\problem{1} Даны окружность $\omega$ и фиксированная точка $A$ вне окружности. Через точку $A$ проводятся окружности $S$, которые касаются окружности $\omega$ в точке $B$. Касательные, проведенные в точках $A$ и $B$ к окружности $S$, пересекаются в точке $M$. Докажите, что все такие точки $M$ лежат на одной прямой.

\problem{2} В остроугольном треугольнике $A B C$ угол $C$ больше угла $A$. $M$ - середина стороны $A C$. $A_1$ и $C_1$ - основания высот, проведенных из вершин $A$ и $C$ соответственно. $A_0$ и $C_0$ - середины отрезков $M A_1$ и $M C_1$. Прямая $A_0 C_0$ пересекает прямую, проходящую через $B$ параллельно $A C$, в точке $T$. Докажите, что $T B=T M$.

\problem{3} Вписанная окружность с центром $I$ касается сторон $A B, B C, A C$ треугольника $A B C$ в точках $C_0, A_0, B_0$. Прямая, перпендикулярная $B B_0$ и проходящая через $B_0$, пересекает $A_0 C_0$ в точке $B_1$. Докажите, что середина отрезка $B B_1$ лежит на прямой $A C$.

\problem{4} Из вершины $C$ треугольника $A B C$ проведены касательные $C X, C Y$ к окружности, проходящей через середины сторон треугольника. Докажите, что прямые $X Y$, $A B$ и касательная в точке $C$ к окружности, описанной около треугольника $A B C$, пересекаются в одной точке.

\problem{5}
Периметр треугольника $A B C$ равен 4. На лучах $A B$ и $A C$ отмечены точки $X$ и $Y$ так, что $A X=A Y=1$. Отрезки $B C$ и $X Y$ пересекаются в точке $M$. Докажите, что периметр одного из треугольников $A B M$ или $A C M$ равен 2 .

\problem{6}
Вписанная окружность треугольника $A B C$ касается сторон $A B$ и $A C$ в точках $Z$ и $Y$ соответственно. Прямые $B Y$ и $C Z$ пересекаются в точке $G$. Точки $R$ и $S$ выбираются так, что четырехугольники $B C Y R$ и $B C S Z-$ параллелограммы. Докажите, что $G R=G S$.

